\subsection*{The received gospel opens a hearing people}

The central relation is not simply ``deafness becomes hearing.'' It is that a
defined word is received before speech becomes right. Paul recalls a gospel
which he first delivered because he had received it: Christ's death for sins,
burial, Resurrection, and appearances according to the Scriptures. Mark then
shows Christ opening a man's ears and releasing the bond of his tongue. The
Alleluia calls the people to jubilant sound but comes from a psalm whose larger
movement demands covenant hearing. The Collect places every such reception
under mercy surpassing the petitioner's merits and requests.

The direct witnesses preserve the order. Chrysostom reads Paul's
``delivered'' and ``received'' as the language of transmission rather than
private speculation. Gregory the Great reads Mark's fingers as gifts of the
Spirit and the saliva as wisdom from divine speech, ending with the demand
that one do before teaching. Bede develops opened ears toward catechetical
confession. Augustine's treatment of the command to silence refuses to make
verbal abundance the measure of obedience: those charged to preach should
learn zeal even from witnesses who were forbidden, but the narrative itself
still distinguishes healed speech from authorized proclamation.

This relation corrects two opposite errors. Christian faith is not bare
information received without embodied encounter: Mark's Christ touches,
looks, groans, and speaks. Nor is religious expression authenticated by its
intensity: Paul gives the confession determinate apostolic content, and the
Gospel's secrecy command interrupts publicity. Hearing is opened toward a
word already given; speech becomes ``right'' when it answers that word in a
life shaped by it.

\subsection*{Grace fruitful without becoming possession}

Four moments define a movement of grace without ownership. The Collect asks
for mercy beyond merit and request. The Epistle ends with grace in Paul not
having become void. The Secret calls the Church's offered service a gift but
asks God to make it acceptable and helpful. The Communion takes substance and
firstfruits into the honor of God. Each human act is real---supplication,
holding fast, preaching, offering, giving---and each remains answerable to a
giver and an end outside itself.

The exact Pauline boundary keeps the formulation honest. The lection does not
include ``I labored more abundantly'' or ``the grace of God with me.'' The
full canonical verse supports Augustine's, Gregory's, and Aquinas's accounts
of grace preceding and cooperating with human freedom, but those accounts may
not be attributed to the fragment alone. Within the appointed text, Paul
already supplies enough for a narrower conclusion: the former persecutor's
apostolic identity is grace's work, and grace has borne fruit rather than
remaining empty.

Bede's exposition of Proverbs~3 adds the needed moral specificity. One honors
God from goods honestly held, gives without seeking applause, and credits
every good work to grace. Salvian makes the economic claim sharper by treating
possessions as first received from God and therefore owed back in devotion and
justice. The Secret brings that outward motion to the altar, but it still asks
for acceptance and remedy. Grace does not abolish agency; it prevents agency,
office, goods, or sacrifice from closing into a private title.

\subsection*{The body is healed toward resurrection}

The formulary's bodily vocabulary is unusually dense: God gives strength;
flesh flourishes again; ears and tongue are restored; the singer is upheld and
healed; fragility receives help; barns and presses overflow; mind and body are
saved. None of these images by itself proves an eschatological design. Their
canonical and reception histories nevertheless establish a coherent bodily
horizon.

Irenaeus and Tertullian defend from 1~Corinthians~15 the real flesh of the
Christ who died, was buried, and rose. Augustine and Cassiodorus interpret the
Gradual's flourishing flesh as Resurrection; Theodoret preserves a literal
alternative of Davidic vigor restored. The same Latin and Greek traditions
read Psalm~29's dedication as the renewal of Christ's body, the Church, or
human nature in Resurrection. Mark supplies a bodily healing good in itself,
not merely material for allegory. Finally, the Postcommunion explicitly holds
mind and body inside one salvation and asks for a heavenly remedy whose
fullness is not exhausted in present cure.

Schuster's Eucharistic reading draws these strands toward final resurrection:
contact with Christ's sacramental Body strengthens the embodied Christian for
the body's promised raising. That is later liturgical theology, not an
ancient gloss on every text. It nonetheless articulates why the sequence may
not be spiritualized. Christian salvation opens hearing and speech, forms
conscience, receives material firstfruits, sustains fragile bodies, and hopes
for incorruption. Present illness is not failed faith, and disability is not
culpability; the controlling promise is the risen Christ.

\subsection*{Silence, groaning, and true praise}

Two psalms ask and answer a problem of voice. The Gradual pleads that God not
be silent or depart; in the Gospel, Jesus looks to heaven, groans, and speaks
one effective word. The Offertory answers healing with extolment; the
Alleluia orders festive sound toward God the helper. The sequence refuses both
wordlessness and chatter. Divine action may include a groan and a single
imperative; human thanksgiving may be abundant, but proclamation must remain
true to what has been received.

Gregory says Christ's sigh was exemplary rather than necessary, directing
our prayerful groaning toward heaven. Theophylact offers several possibilities---filial appeal and
compassion for suffering human nature---without claiming to exhaust the
gesture. Saint Lawrence of Brindisi and later papal catecheses read the
heavenward gaze and \emph{Ephphatha} toward grace, communion, and mission.
Pope Leo~XIV adds a narrative caution: Mark's secrecy belongs to a Gospel in
which a mature confession must pass through the entire paschal journey. Thus
the Gradual's fear of silence is not solved by maximum output but by the Lord
whose word opens and whose Pasch determines speech's content.

\subsection*{A holy household of service and remedy}

The Introit's holy place and concordant house might remain a beautiful image
unless the other propers disclose its practices. Paul addresses an actual
church asked to stand in the received gospel. Unnamed bearers in Mark bring a
particular person into Christ's healing intimacy. The Secret
names worship as service and asks that it become assistance for weakness. The
Communion orders substance and firstfruits toward the Lord. The
Postcommunion makes the recipients supplicants again.

Hilary and Cassiodorus connect Psalm~67's household with the one heart and
soul of Acts~4:32. Augustine insists that grace, not merit, builds the holy
place. Saint John Henry Newman receives Mark's bearers as a figure of the
Church's intercession and mission while expressly preserving divine freedom
to act beyond visible structures. These sources support a churchly reading
only when the healed man remains a person rather than a symbol and the common
house remains ordered to care.

The whole integrated movement can therefore be stated without walking the
ten elements in order. A people made by grace receives a bodily Resurrection
gospel; Christ opens persons to hear and confess it; healed and strengthened
persons return praise, service, and firstfruits; the altar and Communion do
not close the account but make the household ask again for accepted worship
and whole-person remedy. Divine gift becomes communion precisely by refusing
to become private possession.
